\section{Open questions}
The replication verified MLAE's core scaling claim in noiseless simulation.
Five substantive questions the paper (and this replication) leave open:

\begin{enumerate}
  \item \textbf{Noise robustness.} How does the near-Heisenberg EIS scaling
        degrade under realistic gate/measurement noise on current
        superconducting or trapped-ion hardware?  This replication used
        noiseless \texttt{qiskit-aer} only.  MLAE's practical selling point
        for near-term devices is that its shallower circuits (no QFT
        ancilla register) are more noise-robust than canonical QAE, but
        the noise-vs-slope tradeoff was not measured at NISQ error rates.
        \emph{Next:} rerun the $a=1/48$ sweep with a depolarizing noise
        model calibrated to IBM Falcon/Heron error rates, sweeping $p$;
        report the slope surface, then repeat on IBM Quantum's free tier
        for $M\le 5$ EIS.

  \item \textbf{Head-to-head vs canonical QAE.} Does MLAE's advantage survive
        when compared against QPE-based amplitude estimation at matched
        \emph{gate-count} (not matched $N_q$)?  This replication only
        measured RMSE vs oracle-call count; the ancilla/depth savings claim
        was not exercised.  \emph{Next:} implement Brassard--H\o yer--Mosca--Tapp
        QAE with $n_{\text{eval}} \in \{3,\dots,7\}$; overlay against
        MLAE-EIS on RMSE-vs-2Q-gate-count and RMSE-vs-depth axes.

  \item \textbf{Comparison to Iterative QAE (Grinko et al.).} How does MLAE
        compare to the modern successor IAE in shot efficiency at fixed
        target accuracy?  IAE post-dates this paper and claims comparable
        efficiency without the MLE local-minimum risk.  \emph{Next:} use
        Qiskit \texttt{IterativeAmplitudeEstimation} with
        $\varepsilon_{\text{target}}$ swept in $\{10^{-2},\dots,5\!\times\!10^{-4}\}$;
        compare calls-to-target for MLAE-EIS vs IAE at $a=1/48$ and $a=0.3$.

  \item \textbf{MLE branch-identifiability failure mode.} How badly does the
        MLE misidentify the branch when $\theta_a$ is near $\pi/4$ or when
        $N_{\text{shot}}$ per round is very small?  The likelihood surface
        $\sin^{2}((2m_k+1)\theta_a)$ is multimodal; neither the paper nor
        this replication quantifies wrong-branch miss-rate as a function
        of $(a, N_{\text{shot}})$.  \emph{Next:} sweep $a\in\{0.01,\dots,0.49\}$
        and $N_{\text{shot}}\in\{5,10,25,100\}$ at EIS $M=6$, 500 trials
        each; publish a wrong-mode heatmap.

  \item \textbf{Real-application end-to-end (quant finance).} Does MLAE
        deliver its scaling advantage on a specific Monte-Carlo task
        (European option pricing) end-to-end, when $A$ itself is a
        non-trivial state-prep circuit whose depth may dominate small
        $Q^{m_k}$?  \emph{Next:} use Qiskit Finance's
        \texttt{EuropeanCallPricing} as $A$, run MLAE-EIS with $M\in\{3,\dots,7\}$;
        report RMSE(price) vs $N_q$ and identify the $M$ threshold beyond
        which the EIS advantage exceeds $A$'s depth cost.
\end{enumerate}
